\documentclass[12pt]{article}

\usepackage{amsmath}
\usepackage{amssymb}
\usepackage{physics}
\usepackage{hyperref}

\title{Foundations of Affine Space}
\author{}
\date{}

\begin{document}

\maketitle

\section{Introduction}
% ============================================================

\subsection{Notation and conventions}

Before beginning the mathematical development, we establish the notation
that will be used throughout this essay.

Points and vectors are distinct kinds of objects: points belong to
\(\mathcal A\), whereas vectors belong to \(V\).

\begin{center}
\begin{tabular}{lll}
Symbol & Meaning & Convention or remarks \\
\hline
\(\mathcal A\) & Ambient affine space & Its elements are points \(x,y,z\in\mathcal A\). \\
\(V\) & Vector space associated with \(\mathcal A\) & Its elements are bold lowercase vectors \(\mathbf u,\mathbf v,\mathbf w\in V\). \\
\(U,W\) & Linear subspaces of \(V\) & Typical elements are \(\mathbf u\in U\) and \(\mathbf w\in W\). \\
\(\{\mathbf v_i\}_{i=1}^{n}\) & Basis of \(V\) & \(n=\dim\mathcal A=\dim V\). \\
\(\{\mathbf u_i\}_{i=1}^{r}\) & Basis of \(U\) & \(r=\dim U\). \\
\(\{\mathbf w_i\}_{i=1}^{s}\) & Basis of \(W\) & \(s=\dim W\). \\
\(S(x,U)\) & Affine subspace of \(\mathcal A\) through \(x\) with direction space \(U\) & \(x\in\mathcal A\) and \(U\) linear subspace of \(V\). \\
\end{tabular}
\end{center}

The vector \(\overrightarrow{xy}\in V\) denotes the unique displacement
from the point \(x\) to the point \(y\).

Once an affine coordinate system has been chosen, point coordinates and
vector components are distinguished as shown below.

\begin{center}
\begin{tabular}{lll}
Symbol & Meaning & Convention or remarks \\
\hline
\(\overrightarrow{xy}\) & Unique displacement from \(x\) to \(y\) & \(\overrightarrow{xy}\in V\). \\
\(\mathbf d\) & Explicitly named displacement & \(\mathbf d:=\overrightarrow{xy}\). \\
\(\mathbf d_{ij}\) & Indexed displacement & \(\mathbf d_{ij}:=\overrightarrow{x_i x_j}\). \\
\((x^1,\ldots,x^n)\) & Coordinates of \(x\in\mathcal A\) & Defined after choosing an affine coordinate system. \\
\(v^i,u^i,w^i\) & Components of vectors & \(\mathbf v=v^i\mathbf v_i\), \(\mathbf u=u^i\mathbf u_i\), and \(\mathbf w=w^i\mathbf w_i\). \\
\(i,j\) & Basis and coordinate indices & Their range is determined by the relevant basis: \(1,\ldots,n\), \(1,\ldots,r\), or \(1,\ldots,s\). \\
\(v^i\mathbf v_i\) & Einstein summation convention & An index repeated once upstairs and once downstairs is summed over its range. The same convention applies to \(u^i\mathbf u_i\) and \(w^i\mathbf w_i\). \\
\end{tabular}
\end{center}

Additional symbols will be introduced as the corresponding mathematical
objects are needed.

\section{Affine Spaces}

\begin{definition}[Affine space]
An \(n\)-dimensional affine space over \(\mathbb{R}\) consists of:
\begin{itemize}
    \item a vector space \(V\);
    \item a set \(\mathcal A\) of points.
\end{itemize}
\end{definition}

To each vector \(\mathbf{v} \in V\) is associated a mapping
\[
\mathbf{v} \colon \mathcal A \longrightarrow \mathcal A,
\]
\[
\mathbf{v}(x) = y,
\]
where \(x,y \in \mathcal A\), with the following properties:

\begin{enumerate}
    \item
    \(\mathbf{0}(x)=x\), where \(\mathbf{0}\) is the null vector.

    \item
    For all \(x,y\in \mathcal A\), there exists a unique \(\mathbf{v}\in V\)
    such that \(\mathbf{v}(x)=y\), and we denote
    \(\mathbf{v}=\overrightarrow{xy}\).

    \item
    If \(\mathbf{v},\mathbf{w}\in V\), the composition of the action
    of two vectors is defined formally as the sum of two vectors:
    
    \[(\mathbf{v}\circ\mathbf{w})(x)
    =\mathbf{v}\bigl(\mathbf{w}(x)\bigr)
    \coloneqq(\mathbf{v}+\mathbf{w})(x).\]

\end{enumerate}

Such a composition must be commutative, as the sum of vectors is:
\[
\mathbf{v}\circ\mathbf{w}
=
\mathbf{v}+\mathbf{w}
=
\mathbf{w}+\mathbf{v}
=
\mathbf{w}\circ\mathbf{v}.
\] 


\begin{definition}[Translation]
To any given vector \(\mathbf{v}\in V\) is associated a mapping
\[
T_{\mathbf{v}}\colon \mathcal A\longrightarrow \mathcal A
\]
such that
\[
T_{\mathbf{v}}(x)=\mathbf{v}(x).
\]
Such a mapping is called a translation.
\end{definition}

\begin{proposition}
For every \(\mathbf{v}\in V\), the translation
\[
T_{\mathbf{v}}\colon \mathcal A\longrightarrow \mathcal A
\]
is one-to-one and onto.
\end{proposition}

\begin{proof}
Let us prove that \(T_{\mathbf{v}}\) is one-to-one. Assume that
\(z,y\in \mathcal A\) and
\[
T_{\mathbf{v}}(z)=T_{\mathbf{v}}(y).
\]
Therefore,
\[
\mathbf{v}(z)=\mathbf{v}(y).
\]
Applying \(-\mathbf{v}\) to both sides, we obtain
\[
(-\mathbf{v})\bigl(\mathbf{v}(z)\bigr)
=
(-\mathbf{v})\bigl(\mathbf{v}(y)\bigr).
\]
Hence,
\[
(-\mathbf{v}+\mathbf{v})(z)
=
(-\mathbf{v}+\mathbf{v})(y),
\]
and therefore
\[
\mathbf{0}(z)=\mathbf{0}(y).
\]
It follows that
\[
z=y.
\]

Let us now prove that \(T_{\mathbf{v}}\) is onto. For every \(y\in \mathcal A\), we
need to prove that there exists \(x\in \mathcal A\) such that
\[
\mathbf{v}(x)=y.
\]
We simply take
\[
x=(-\mathbf{v})(y).
\]
Then
\[
\mathbf{v}(x)
=
\mathbf{v}\bigl((-\mathbf{v})(y)\bigr)
=
(\mathbf{v}-\mathbf{v})(y)
=
\mathbf{0}(y)
=
y.
\]
\end{proof}

\begin{definition}[Permutation group]
Let \(Y\) be a set. The set
\[
\operatorname{Sym}(Y)
=
\left\{
f\colon Y\longrightarrow Y
\;\middle|\;
f \text{ is bijective}
\right\}
\]
is called the permutation group of \(Y\). Its group operation is the
composition of mappings.
\end{definition}

\begin{proposition}
The mapping
\[
\rho\colon V\longrightarrow\operatorname{Sym}(\mathcal A),
\qquad
\rho(\mathbf{v})=T_{\mathbf{v}},
\]
is an injective homomorphism from the additive group of \(V\) into the
permutation group \(\operatorname{Sym}(\mathcal A)\).
\end{proposition}

\begin{proof}
We first prove that \(\rho\) preserves the group operation. Let
\(\mathbf{u},\mathbf{v}\in V\). For every \(x\in \mathcal A\),
\[
\rho(\mathbf{u}+\mathbf{v})(x)
=
T_{\mathbf{u}+\mathbf{v}}(x)
=
(\mathbf{u}+\mathbf{v})(x).
\]
By the definition of the action of the sum of two vectors,
\[
(\mathbf{u}+\mathbf{v})(x)
=
\mathbf{u}\bigl(\mathbf{v}(x)\bigr)
=
\bigl(T_{\mathbf{u}}\circ T_{\mathbf{v}}\bigr)(x).
\]
Therefore,
\[
\rho(\mathbf{u}+\mathbf{v})
=
\rho(\mathbf{u})\circ\rho(\mathbf{v}),
\]
and hence \(\rho\) is a homomorphism.

We now prove that \(\rho\) is injective. Assume that
\[
\rho(\mathbf{u})=\rho(\mathbf{v}).
\]
Then
\[
T_{\mathbf{u}}(x)=T_{\mathbf{v}}(x)
\]
for every \(x\in \mathcal A\), and therefore
\[
\mathbf{u}(x)=\mathbf{v}(x).
\]
Applying \(-\mathbf{v}\) to both sides gives
\[
(-\mathbf{v})\bigl(\mathbf{u}(x)\bigr)
=
(-\mathbf{v})\bigl(\mathbf{v}(x)\bigr).
\]
Thus,
\[
(-\mathbf{v}+\mathbf{u})(x)
=
(-\mathbf{v}+\mathbf{v})(x)
=
\mathbf{0}(x)
=
x.
\]
By the uniqueness of the vector that maps \(x\) onto itself,
\[
-\mathbf{v}+\mathbf{u}=\mathbf{0},
\]
and hence
\[
\mathbf{u}=\mathbf{v}.
\]
Therefore, \(\rho\) is injective.
\end{proof}


\end{document}
